\documentclass[11pt]{article}
\usepackage[margin=0.9in]{geometry}
\usepackage{amsmath,amssymb,amsfonts}
\newcommand{\Lam}{\Lambda}
\newcommand{\R}{\mathbb{R}}
\pagestyle{empty}
\begin{document}
\begin{center}{\Large\textbf{LGM-SSP — latent linear Hawkes superposition, closed-form rate pin}}\end{center}
\vspace{0.4em}

\noindent\textbf{Idea.}\; The rate is a \emph{linear} read-out of a latent state-space
(latent linear Hawkes, LLH) \textbf{superposition}, and its mean is calibrated to $R_{\mathrm{target}}$
in \textbf{closed form} (no EMA / running estimate). One field gives both the rate (its sum) and the
marks (its normalized composition).

\medskip
\noindent\textbf{Latent superposition.}\; $x(t)\in\R^{P}$ is a superposition of $P$ decaying modes;
between events it decays, and each event of mark $j$ adds a non-negative impulse $B_j$:
\[
x(t)=\sum_{t_i<t} B_{k_i}\,e^{-\delta\,(t-t_i)}\qquad(\delta_p>0,\ B\ge0,\ \text{elementwise per mode}).
\]
\noindent\textbf{Per-mark intensities (multivariate linear Hawkes), rate, marks.}
\[
\boxed{\;\lambda_k(t)=\nu_k+(W\,x(t))_k\;}\quad(W\ge0,\ \nu\ge0),\qquad
\Lam(t)=\sum_{k=1}^{K}\lambda_k(t),\qquad p(k\mid t)=\frac{\lambda_k(t)}{\Lam(t)} .
\]

\medskip
\noindent\textbf{Closed-form calibration (the pin) --- no EMA.}\;
With the branching matrix and ratio
\[
G_{kj}=\sum_{p=1}^{P}\frac{W_{kp}\,B_{jp}}{\delta_p},\qquad
n=\rho(G)\ \text{(spectral radius, projected to }n\le n_{\mathrm{cap}}<1),
\]
the linear-Hawkes stationary mean is the fixed point $\mathbb{E}[\lambda]=(I-G)^{-1}\nu$. Rescaling the
baseline once,
\[
\boxed{\;\nu=\nu_{\mathrm{raw}}\cdot\frac{R_{\mathrm{target}}}{\mathbf 1^{\top}(I-G)^{-1}\nu_{\mathrm{raw}}}
\quad\Longrightarrow\quad
\mathbb{E}\Big[\textstyle\sum_k\lambda_k\Big]=\mathbf 1^{\top}(I-G)^{-1}\nu=R_{\mathrm{target}}\;}
\]
pins the total stationary mean \emph{exactly} and analytically (verified: $=40.600$ for
$R_{\mathrm{target}}=40.6$). $n<1$ is a genuine subcriticality certificate.

\medskip
\noindent\textbf{In words.}\; A multivariate (cross-exciting) linear Hawkes built on a latent decaying
superposition: the total rate is $\sum_k\lambda_k$, the marks are $\lambda_k/\sum_j\lambda_j$, and the
baseline $\nu$ is set in closed form (one $(I-G)^{-1}$ solve) so the mean event rate is exactly
$R_{\mathrm{target}}$. This is LGM's exact-pin philosophy ($\mu_0=R_{\mathrm{target}}(1-n)$) generalized
from a scalar ground rate to the full latent-linear-Hawkes superposition --- replacing the earlier
EMA rescale with an analytic stationary-mean pin, and recovering the $n<1$ certificate.

\medskip
\noindent\textbf{Likelihood.}\;
$\sum_i\log\lambda_{k_i}(t_i)-\int_0^T\Lam(s)\,ds$; the compensator is closed-form per mode
($\int e^{-\delta_p\tau}$), so no Monte-Carlo is needed for the rate. Linear intensity is kept
non-negative by $W,\nu,B\ge0$.
\end{document}
